\chapter{Directed Acyclic Graph Convolutional Networks: Learning Causal Signals}
\label{chap:rey_dagcn}

% =============================================================================
\section{Why Standard GNNs Cannot Learn on DAGs}
\label{sec:dagcn_motivation}
% =============================================================================

Directed acyclic graphs (DAGs) are one of the most important families of graphs
in applied science. A gene regulatory network is a DAG: Gene A activates Gene
B, and Gene B activates Gene C, but the regulation flows in one direction only
--- Gene C cannot retroactively regulate Gene A. A Bayesian network is a DAG:
each node is a random variable, and directed edges encode conditional
dependencies from causes to effects. A neural architecture search space is a
DAG: layers are composed in topological order. A workflow scheduler is a DAG:
task dependencies form an acyclic precedence relation. In each of these domains,
the DAG structure encodes something fundamental about the system --- the
\emph{causal direction} of influence.

It is therefore natural to ask: can we learn from signals defined on DAGs using
Graph Neural Networks (GNNs)? Can we build a GNN that respects the directional,
causal structure of a DAG?

The answer, for standard GNNs, is: surprisingly poorly. Rey, Ajorlou, and
Mateos~\cite{rey2025dagcn} identify two fundamental failures that arise when
standard spectral or message-passing GNNs are applied to DAGs. Understanding
these failures is the starting point for their \emph{Directed Acyclic Graph
Convolutional Network} (DAGCN), which we study in this chapter.

\subsection{The Standard Filterbank GNN}

The standard spectral GNN, often called a \textbf{graph filterbank convolutional
neural network} (FB-GCNN), applies repeated polynomial filtering of the input
signal matrix $\mathbf{X} \in \mathbb{R}^{N \times F}$ using powers of the
adjacency matrix $\mathbf{A}$:

\begin{equation}
    \mathbf{X}^{(\ell+1)} = \sigma\!\left(
        \sum_{r=0}^{R} \mathbf{A}^r \, \mathbf{X}^{(\ell)} \, \boldsymbol{\Theta}_r^{(\ell)}
    \right)
    \label{eq:fb_gcnn}
\end{equation}

where $\boldsymbol{\Theta}_r^{(\ell)} \in \mathbb{R}^{F_\ell \times F_{\ell+1}}$
are learnable weight matrices and $\sigma$ is a pointwise nonlinearity such as
ReLU. The term $\mathbf{A}^r \mathbf{X}^{(\ell)}$ computes the $r$-hop
neighbourhood aggregation: each node gathers information from nodes up to $r$
steps away.

This architecture has been enormously successful on undirected and general
directed graphs. On DAGs, however, it encounters two structural problems.

\subsection{Failure Mode 1: The Nilpotent Spectrum}

\begin{conceptbox}{Nilpotency of the DAG Adjacency Matrix}
\textbf{Definition.} A matrix $\mathbf{A}$ is \textbf{nilpotent} if there
exists an integer $n \ge 1$ such that $\mathbf{A}^n = \mathbf{0}$.

\medskip
\textbf{Claim.} For any DAG on $N$ nodes with adjacency matrix
$\mathbf{A} \in \{0,1\}^{N \times N}$, we have $\mathbf{A}^N = \mathbf{0}$.

\medskip
\textbf{Why.} A DAG has no cycles. Order the nodes in topological order:
label node 1 as the first node in the topological sequence, node 2 as the
second, and so on. In this ordering, every edge $(i, j) \in E$ satisfies
$i > j$ (edges only go from lower-indexed to higher-indexed nodes in
topological order). The matrix $\mathbf{A}$ is therefore \emph{strictly
lower-triangular} in this ordering. For any strictly lower-triangular matrix,
multiplying by $\mathbf{A}$ shifts all non-zero entries strictly downward
along the diagonal. After $N$ such multiplications, all entries have been
shifted off the bottom of the matrix and $\mathbf{A}^N = \mathbf{0}$.

\medskip
\textbf{Consequence for eigenvalues.}
The characteristic polynomial of a strictly lower-triangular matrix is
$\det(\mathbf{A} - \lambda \mathbf{I}) = (-\lambda)^N$, so all $N$
eigenvalues of $\mathbf{A}$ are \textbf{zero}. The spectral decomposition
$\mathbf{A} = \mathbf{U} \boldsymbol{\Lambda} \mathbf{U}^{-1}$ (with
$\boldsymbol{\Lambda} = \mathbf{0}$) is degenerate --- all eigenvectors
correspond to the same eigenvalue zero, and the matrix is not diagonalisable
in general.
\end{conceptbox}

This single fact collapses the theoretical foundation of spectral GNNs on DAGs.
Spectral graph convolution is defined via the graph's eigenbasis: the graph
Fourier transform decomposes a signal into components corresponding to
different graph frequencies (eigenvalues). If all eigenvalues are zero, all
``frequencies'' are identical, and the Fourier transform provides no
discriminative information. Every node looks the same from the spectral
perspective, regardless of its position in the DAG.

\begin{examplebox}{Nilpotency in the Gene Network}
Consider our running four-node gene regulatory network:
\begin{itemize}[noitemsep]
    \item Nodes: Gene A, Gene B, Gene C, Gene D.
    \item Edges (regulatory): A $\to$ C, B $\to$ C, C $\to$ D.
\end{itemize}
In topological order A, B, C, D (labelled 1, 2, 3, 4), the adjacency matrix is:
\[
    \mathbf{A} = \begin{pmatrix} 0 & 0 & 0 & 0 \\ 0 & 0 & 0 & 0 \\ 1 & 1 & 0 & 0 \\ 0 & 0 & 1 & 0 \end{pmatrix}
\]
where $A_{ij} \neq 0$ means there is a directed edge \emph{from} $j$ \emph{to}
$i$ (i.e., gene $j$ regulates gene $i$). Compute:
\[
    \mathbf{A}^2 = \begin{pmatrix} 0&0&0&0\\0&0&0&0\\0&0&0&0\\1&1&0&0 \end{pmatrix}, \quad
    \mathbf{A}^3 = \begin{pmatrix} 0&0&0&0\\0&0&0&0\\0&0&0&0\\0&0&0&0 \end{pmatrix} = \mathbf{0}
\]
After just 3 multiplications (fewer than $N = 4$), the matrix is zero.
All eigenvalues of $\mathbf{A}$ are zero, confirming nilpotency. A spectral
GNN using polynomial filters $\sum_r h_r \lambda^r$ evaluates to $h_0 \mathbf{I}$
on every node --- it can only apply a global scalar shift, learning nothing about
the DAG structure.
\end{examplebox}

\subsection{Failure Mode 2: Ignoring the Partial Order}

Even if we set aside spectral methods and consider spatial (message-passing)
GNNs, a second and more fundamental problem remains. Standard message passing
aggregates information from all neighbours, regardless of the direction of
edges. On a DAG, however, direction \emph{is} the structure. An edge $j \to i$
means ``$j$ causally influences $i$'', not the reverse.

\begin{keyinsight}{Direction Encodes Causality --- Ignoring It Is Wrong}
In a gene regulatory network, if Gene A regulates Gene C, then Gene A's
expression level is a \emph{cause} of Gene C's expression level. A GNN that
aggregates information symmetrically --- allowing Gene C's expression to
influence Gene A's representation just as much as Gene A's expression influences
Gene C's --- is learning a model that is causally inverted with respect to the
actual biology.

The problem is not merely a loss of information. It is the introduction of
\emph{incorrect} information: backward message passing from descendants to
ancestors implies a causal influence that does not exist. For downstream tasks
that depend on causal reasoning (such as source detection, causal inference, or
circuit simulation), this contamination of representations is not a minor error
--- it is a fundamental mis-specification.
\end{keyinsight}

\paragraph{Worked example: phantom backward flow.}
Return to the gene network. Gene D is purely downstream: it is regulated by C,
which is regulated by A and B. Gene D has no regulatory influence on any other
gene. In a standard undirected or bidirectional GNN, Gene D's expression signal
would be propagated back to Gene C during message passing, and from C back to
Genes A and B. After two message-passing steps, the representation of Gene A
would contain contributions from Gene D. This is biologically nonsensical: Gene
D has no regulatory effect on Gene A in this network. The standard GNN has
created a phantom causal pathway $D \to C \to A$ that does not exist.

\begin{analogybox}{A One-Way Street That GNNs Drive the Wrong Way}
Imagine a city whose streets form a perfect one-way system: you can drive from
the city centre outward, but never inward. A delivery driver using a map that
ignores one-way signs would plan routes that are physically illegal and
impossible. Standard GNNs are like that map: they see connections but not
their direction, and they plan impossible routes through the network as a result.

The DAGCN paper builds a map that is aware of one-way streets --- one that
respects the causal direction of every edge and builds representations
accordingly. The key insight is that we need a new notion of ``graph frequency''
that respects direction: not the eigenvectors of $\mathbf{A}$ (which are
degenerate), but the columns of a structure that captures the full
\emph{ancestry} relations in the DAG.
\end{analogybox}

% =============================================================================
\section{DAGs, Partial Orders, and What ``Comes Before'' Means}
\label{sec:dagcn_partial_order}
% =============================================================================

To build a proper theory of signal processing on DAGs, we first need a precise
language for the \emph{precedence structure} that a DAG encodes. This section
develops that language carefully, with explicit attention to the running gene
network example.

\subsection{From Directed Edges to a Partial Order}

\begin{mathbox}{Definition: Partial Order Induced by a DAG}
Let $D = (\mathcal{V}, \mathcal{E})$ be a DAG with $|\mathcal{V}| = N$ nodes.
We say that node $j$ \textbf{precedes} node $i$ (written $j \prec i$, or
equivalently $j < i$) if there exists a directed path from $j$ to $i$ in $D$.

The relation $\prec$ satisfies the three axioms of a \textbf{partial order}:
\begin{enumerate}[noitemsep]
    \item \textbf{Irreflexivity:} No node precedes itself ($j \not\prec j$),
          because a directed path from $j$ to $j$ would form a cycle, contradicting
          the definition of a DAG.
    \item \textbf{Asymmetry:} If $j \prec i$ then $i \not\prec j$, because a
          directed path from $j$ to $i$ and another from $i$ to $j$ would together
          form a cycle.
    \item \textbf{Transitivity:} If $j \prec k$ and $k \prec i$ then $j \prec i$,
          because the two directed paths can be concatenated.
\end{enumerate}

We write $j \le i$ to mean either $j = i$ or $j \prec i$ (reflexive closure).
In the matrix convention of Rey et al., $A_{ij} \neq 0$ means there is a direct
edge from $j$ to $i$, so $j \prec i$ directly.
\end{mathbox}

Unlike a \emph{total} order (where every pair of elements is comparable), a
partial order allows pairs of nodes that are \emph{incomparable}: neither
$j \prec i$ nor $i \prec j$. In the gene network, Genes A and B are
incomparable --- neither regulates the other.

\subsection{Topological Ordering}

A \textbf{topological ordering} of a DAG assigns a label $\ell: \mathcal{V}
\to \{1, \dots, N\}$ to each node such that if $j \prec i$, then
$\ell(j) < \ell(i)$. Every finite DAG has at least one topological ordering
(and generally many). The partial order is consistent with any topological
ordering, but it is richer: it records which nodes are \emph{related} by the
precedence relation, not just the numeric labels.

\begin{examplebox}{Partial Order in the Gene Regulatory Network}
Our four-node network: edges A $\to$ C, B $\to$ C, C $\to$ D. Recall the
convention: $A_{ij} \neq 0$ means edge \emph{from} $j$ \emph{to} $i$.

We enumerate the precedence relationships:
\begin{center}
\begin{tabular}{lll}
\toprule
\textbf{Pair} & \textbf{Directed path exists?} & \textbf{Relation} \\
\midrule
A, C & A $\to$ C (direct edge) & $A \prec C$ \\
B, C & B $\to$ C (direct edge) & $B \prec C$ \\
C, D & C $\to$ D (direct edge) & $C \prec D$ \\
A, D & A $\to$ C $\to$ D (length-2 path) & $A \prec D$ \\
B, D & B $\to$ C $\to$ D (length-2 path) & $B \prec D$ \\
A, B & No path in either direction & Incomparable \\
\bottomrule
\end{tabular}
\end{center}

A valid topological ordering is: A=1, B=2, C=3, D=4. Another valid ordering
is B=1, A=2, C=3, D=4. Both are consistent with the partial order.

The \textbf{predecessor sets} (all nodes that come before a given node,
including itself) are:
\begin{align*}
    \text{pred}(A) &= \{A\} \\
    \text{pred}(B) &= \{B\} \\
    \text{pred}(C) &= \{A, B, C\} \\
    \text{pred}(D) &= \{A, B, C, D\}
\end{align*}

Genes A and B are \emph{root genes}: no gene precedes them. Gene D is the
\emph{terminal gene}: every gene is a predecessor of D (D comes last in every
valid topological ordering).
\end{examplebox}

\begin{figure}[H]
\centering
\begin{tikzpicture}[
    gnode/.style={circle, draw=green!50!black, fill=green!8!white, thick,
                  minimum size=1.1cm, font=\bfseries\sffamily},
    arr/.style={-{Stealth[length=3mm]}, thick, green!40!black},
    lbl/.style={font=\scriptsize\sffamily, fill=white, inner sep=1pt}
]
    \node[gnode] (A) at (0,  0)   {A};
    \node[gnode] (B) at (2,  0)   {B};
    \node[gnode] (C) at (1, -2)   {C};
    \node[gnode] (D) at (1, -4)   {D};

    \draw[arr] (A) -- node[lbl, left]{\tiny reg.} (C);
    \draw[arr] (B) -- node[lbl, right]{\tiny reg.} (C);
    \draw[arr] (C) -- node[lbl, right]{\tiny reg.} (D);

    \node[font=\scriptsize\sffamily\itshape, text=gray] at (-1.2, 0)  {roots};
    \node[font=\scriptsize\sffamily\itshape, text=gray] at (-1.2,-4)  {terminal};
\end{tikzpicture}
\caption{The four-node \textbf{gene regulatory DAG} used as a running example
throughout this chapter. Nodes A and B are root genes (no regulators). Gene C
is regulated by both A and B. Gene D is regulated by C alone. Arrows point
from regulator to target: $A_{ij} \neq 0$ means $j$ regulates $i$.}
\label{fig:gene_dag}
\end{figure}

\begin{analogybox}{The Partial Order as a Bureaucratic Org Chart}
Imagine a company hierarchy. The CEO is at the top. Instructions flow
\emph{downward}: the CEO directs the VP, the VP directs the Manager, the
Manager directs the Engineer. An instruction from the Engineer never flows
upward to the CEO. Two engineers in separate divisions are \emph{incomparable}:
neither directs the other.

The DAG partial order is exactly this org chart. ``$j \prec i$'' means ``$j$
gives instructions that eventually reach $i$'' --- $j$ is above $i$ in the
hierarchy. A GNN that allows backward message passing is like a model of
corporate communication in which engineers can send directives to the CEO.
This is not just inefficient; it is a model of a completely different
organisation.
\end{analogybox}

% =============================================================================
\section{The Transitive Closure --- The DAG's Memory}
\label{sec:dagcn_transitive_closure}
% =============================================================================

The adjacency matrix $\mathbf{A}$ records only \emph{direct} edges: one-step
connections. But on a DAG, what matters for signal propagation is the
\emph{entire ancestry}: all paths of all lengths from every ancestor to every
descendant. The object that captures this is the \textbf{weighted transitive
closure} $\mathbf{W}$.

\subsection{Motivation: Why Direct Edges Are Not Enough}

Consider Gene D in our running example. Its expression is influenced not just
by its direct regulator C, but by C's regulators A and B as well. If Gene A is
highly expressed, it activates Gene C, which in turn activates Gene D. The
``memory'' of Gene A's activity is encoded in Gene D's expression level through
the indirect pathway A $\to$ C $\to$ D. The adjacency matrix $\mathbf{A}$
records only the C $\to$ D edge; it is blind to the indirect influence. The
transitive closure records all of it.

\subsection{Definition and Computation}

\begin{mathbox}{Weighted Transitive Closure}
Let $\mathbf{A}$ be the adjacency matrix of a DAG, ordered topologically so
that $\mathbf{A}$ is strictly lower-triangular. Define:

\begin{equation}
    \mathbf{W} = (\mathbf{I} - \mathbf{A})^{-1}
    = \mathbf{I} + \mathbf{A} + \mathbf{A}^2 + \cdots + \mathbf{A}^{N-1}
    \label{eq:transitive_closure}
\end{equation}

The infinite series terminates at $\mathbf{A}^{N-1}$ because $\mathbf{A}$ is
nilpotent ($\mathbf{A}^N = \mathbf{0}$). The matrix $(\mathbf{I} - \mathbf{A})$
is always invertible because its determinant equals 1 (all diagonal entries are
1, all entries above the diagonal are 0 in topological order, so
$\det(\mathbf{I} - \mathbf{A}) = \prod_i (\mathbf{I} - \mathbf{A})_{ii} = 1$).

\medskip
\textbf{Interpretation.}
$W_{ij} \neq 0$ if and only if $j = i$ (same node) or there is a directed
path from $j$ to $i$ in the DAG. Specifically:
\begin{itemize}[noitemsep]
    \item $W_{ii} = 1$ for all $i$ (from the identity term $\mathbf{I}$).
    \item $W_{ij} = A_{ij} = 1$ if there is a direct edge from $j$ to $i$.
    \item $W_{ij} = [\mathbf{A}^2]_{ij} + \cdots$ accumulates contributions
          from all indirect paths of length 2 and above.
\end{itemize}
In this way, $\mathbf{W}$ is the ``memory'' of the entire DAG: it encodes
all ancestral relationships, not just immediate parents.
\end{mathbox}

\subsection{Step-by-Step Computation for the Gene Network}

We compute $\mathbf{W}$ for the four-node gene regulatory DAG. Recall, in
topological order A=1, B=2, C=3, D=4:

\[
    \mathbf{A} = \begin{pmatrix} 0&0&0&0\\0&0&0&0\\1&1&0&0\\0&0&1&0 \end{pmatrix}
\]

\begin{examplebox}{Computing $\mathbf{W} = \mathbf{I} + \mathbf{A} + \mathbf{A}^2 + \mathbf{A}^3$}

\textbf{Step 1:} Compute $\mathbf{A}^2$. The $(i,j)$ entry counts paths of
length exactly 2 from $j$ to $i$.
\[
    \mathbf{A}^2 = \mathbf{A} \cdot \mathbf{A} = \begin{pmatrix} 0&0&0&0\\0&0&0&0\\0&0&0&0\\1&1&0&0 \end{pmatrix}
\]
The non-zero entries $[\mathbf{A}^2]_{D,A} = 1$ and $[\mathbf{A}^2]_{D,B} = 1$
correspond to the paths A $\to$ C $\to$ D and B $\to$ C $\to$ D respectively.

\medskip
\textbf{Step 2:} Compute $\mathbf{A}^3 = \mathbf{A} \cdot \mathbf{A}^2$.
\[
    \mathbf{A}^3 = \begin{pmatrix} 0&0&0&0\\0&0&0&0\\0&0&0&0\\0&0&0&0 \end{pmatrix} = \mathbf{0}
\]
As expected, $\mathbf{A}^3 = \mathbf{0}$ (the network has no paths of length 3).

\medskip
\textbf{Step 3:} Sum: $\mathbf{W} = \mathbf{I} + \mathbf{A} + \mathbf{A}^2$.
\[
    \mathbf{W} = \begin{pmatrix}1&0&0&0\\0&1&0&0\\0&0&1&0\\0&0&0&1\end{pmatrix}
    + \begin{pmatrix}0&0&0&0\\0&0&0&0\\1&1&0&0\\0&0&1&0\end{pmatrix}
    + \begin{pmatrix}0&0&0&0\\0&0&0&0\\0&0&0&0\\1&1&0&0\end{pmatrix}
    = \begin{pmatrix}1&0&0&0\\0&1&0&0\\1&1&1&0\\1&1&1&1\end{pmatrix}
\]

Read the entries: $W_{CA} = 1$ (A is an ancestor of C), $W_{CB} = 1$ (B is an
ancestor of C), $W_{DA} = 1$, $W_{DB} = 1$, $W_{DC} = 1$ (A, B, C are all
ancestors of D). The structure of $\mathbf{W}$ perfectly mirrors the DAG's
ancestry relation.
\end{examplebox}

\subsection{The Linear Structural Equation Model}

The transitive closure $\mathbf{W}$ has a beautiful probabilistic
interpretation via the \textbf{Linear Structural Equation Model} (SEM). The
SEM is the standard generative model for signals on DAGs in the causal
inference literature. Each node $i$ has a signal value $x_i$ that depends
linearly on the signals of its parents plus an exogenous (independent) noise
contribution $c_i$:

\begin{equation}
    x_i = \sum_{j: j \prec i} A_{ij} \, x_j + c_i
    \quad \Longleftrightarrow \quad
    \mathbf{x} = \mathbf{A} \mathbf{x} + \mathbf{c}
    \label{eq:sem}
\end{equation}

Solving for $\mathbf{x}$:
\[
    (\mathbf{I} - \mathbf{A})\mathbf{x} = \mathbf{c}
    \quad \Longrightarrow \quad
    \mathbf{x} = (\mathbf{I} - \mathbf{A})^{-1} \mathbf{c} = \mathbf{W} \mathbf{c}
\]

The signal $\mathbf{x}$ at every node is a linear combination of the exogenous
contributions $\mathbf{c}$, mediated by the transitive closure $\mathbf{W}$.

\begin{examplebox}{SEM Interpretation in the Gene Network}
In the gene network, the SEM $\mathbf{x} = \mathbf{W}\mathbf{c}$ gives:
\begin{align*}
    x_A &= c_A \\
    x_B &= c_B \\
    x_C &= c_A + c_B + c_C \\
    x_D &= c_A + c_B + c_C + c_D
\end{align*}
Gene A's expression is purely its own exogenous contribution. Gene C's
expression is the sum of contributions from its two regulators (A and B) plus
its own baseline. Gene D accumulates the contributions of all four genes in the
network --- it carries the full ``memory'' of the regulatory cascade.

\medskip
\textbf{The DAG Fourier Transform.} Rey et al.\ interpret $\mathbf{W}$ as the
\textbf{DAG Fourier basis}: it decomposes the observed signal $\mathbf{x}$ into
the ``pure'' exogenous contributions $\mathbf{c} = \mathbf{W}^{-1}\mathbf{x}$.
Just as the classical Fourier transform decomposes a signal into sinusoidal
frequencies, the DAG inverse transform $\mathbf{c} = \mathbf{W}^{-1}\mathbf{x}$
recovers the ``causal frequencies'' --- how much each node contributes
independently of its ancestors.

Note that $\mathbf{W}^{-1} = (\mathbf{I} - \mathbf{A})$, which gives:
\[
    \mathbf{W}^{-1} = \begin{pmatrix}1&0&0&0\\0&1&0&0\\-1&-1&1&0\\0&0&-1&1\end{pmatrix}
\]
\end{examplebox}

\begin{examplebox}{Numerical Example: Recovering Exogenous Contributions}
Suppose we observe gene expression levels $\mathbf{x} = [2, 3, 7, 9]^\top$
(Genes A, B, C, D respectively). What are the exogenous contributions
$\mathbf{c} = \mathbf{W}^{-1}\mathbf{x}$?

\[
    \mathbf{c} = \begin{pmatrix}1&0&0&0\\0&1&0&0\\-1&-1&1&0\\0&0&-1&1\end{pmatrix}
    \begin{pmatrix}2\\3\\7\\9\end{pmatrix}
    = \begin{pmatrix}2\\3\\7-2-3\\9-7\end{pmatrix}
    = \begin{pmatrix}2\\3\\2\\2\end{pmatrix}
\]

Interpretation: Gene A contributes 2 ``units'' of expression purely from its
own baseline. Gene B contributes 3 units. After removing the inherited
influence of A and B (which contributes $2+3=5$ to Gene C's expression of 7),
Gene C has its own exogenous contribution of 2. Gene D's observed expression
of 9 is Gene C's total expression (7) plus Gene D's own exogenous contribution
of 2. The DAG Fourier transform cleanly separates these inherited and
intrinsic components.
\end{examplebox}

% =============================================================================
\section{Causal Graph Shift Operators --- Projecting onto Causal Channels}
\label{sec:dagcn_gso}
% =============================================================================

With the SEM and the DAG Fourier basis $\mathbf{W}$ established, Rey et al.\
can now define a new family of graph shift operators (GSOs) that replace the
adjacency matrix in the GNN architecture. These \textbf{causal GSOs} are the
technical heart of the paper.

\subsection{Intuition: Causal Channels}

The standard graph shift operator on an undirected graph is the adjacency
matrix $\mathbf{A}$ (or the normalised Laplacian $\mathbf{L}$). Applying
$\mathbf{A}$ to a signal $\mathbf{x}$ at node $i$ produces
$[\mathbf{A}\mathbf{x}]_i = \sum_{j \in \mathcal{N}(i)} x_j$: each node
receives the sum of its neighbours' signals. This is agnostic to direction.

On a DAG, we want something richer: for each node $k$, a shift operator
$\mathbf{S}_k$ that asks the question:

\begin{center}
\emph{``How much of node $i$'s signal is attributable to the causal channel
running through node $k$?''}
\end{center}

Equivalently, $\mathbf{S}_k$ projects $\mathbf{x}$ onto the set of exogenous
contributions $c_j$ for which node $j$ is a \emph{common ancestor} of both
node $i$ and node $k$. The answer gives a ``causal footprint'' of node $k$'s
ancestry as seen from every other node.

\begin{analogybox}{GSOs as Radio Stations}
Think of the DAG as a radio broadcasting network. Each root gene is a
transmitter station. Every node in the DAG receives signals from all transmitter
stations that have a directed path to it. The causal GSO $\mathbf{S}_k$
is a \textbf{frequency selector} tuned to ``the channel of node $k$'': it
passes through only the frequency components (exogenous contributions $c_j$)
that are relevant to node $k$'s causal history, and zeros out all other
frequencies. Different DAGs have different causal channels: no two
non-isomorphic DAGs produce the same set of GSOs.
\end{analogybox}

\subsection{Formal Definition}

\begin{mathbox}{Causal Graph Shift Operator $\mathbf{S}_k$}
For each node $k \in \mathcal{V}$, define the diagonal matrix
$\mathbf{D}_k \in \{0,1\}^{N \times N}$ by:
\[
    [\mathbf{D}_k]_{ii} = \begin{cases} 1 & \text{if } i \le k \; (i = k \text{ or } i \prec k) \\ 0 & \text{otherwise} \end{cases}
\]
In words, $\mathbf{D}_k$ is the diagonal indicator matrix for the set of nodes
that \emph{precede} $k$ (including $k$ itself). Then the \textbf{causal graph
shift operator} for node $k$ is:
\begin{equation}
    \mathbf{S}_k = \mathbf{W} \mathbf{D}_k \mathbf{W}^{-1}
    \label{eq:causal_gso}
\end{equation}

The output $[\mathbf{S}_k \mathbf{x}]_i$ is:
\[
    [\mathbf{S}_k \mathbf{x}]_i = \sum_{\substack{j \le i \\ j \le k}} W_{ij} \, c_j
\]
--- the sum of exogenous contributions $c_j$ over all nodes $j$ that are common
predecessors of both $i$ and $k$, weighted by the ancestry weight $W_{ij}$.
\end{mathbox}

Two properties are immediately noteworthy:

\paragraph{Idempotence.} $\mathbf{S}_k^2 = \mathbf{S}_k$. Applying the
causal shift for node $k$ twice is identical to applying it once. This is
because $\mathbf{D}_k$ is a projection matrix ($\mathbf{D}_k^2 = \mathbf{D}_k$),
and the GSO inherits this property.

\paragraph{Uniqueness.} The full collection $\{\mathbf{S}_k\}_{k \in \mathcal{V}}$
uniquely identifies the DAG (up to isomorphism). No two non-isomorphic DAGs
produce the same set of causal GSOs (Proposition~2 in the paper). This means
the GSOs are expressive enough to distinguish any pair of structurally distinct
DAGs --- an important property for graph learning.

\subsection{Step-by-Step Computation of All Four GSOs}

We now compute all four causal GSOs for the gene regulatory network.

\paragraph{Recall.} $\mathbf{W} = \begin{pmatrix}1&0&0&0\\0&1&0&0\\1&1&1&0\\1&1&1&1\end{pmatrix}$
and $\mathbf{W}^{-1} = \begin{pmatrix}1&0&0&0\\0&1&0&0\\-1&-1&1&0\\0&0&-1&1\end{pmatrix}$.

\begin{examplebox}{Computing $\mathbf{S}_A$ (GSO for Gene A)}

$\text{pred}(A) = \{A\}$, so $\mathbf{D}_A = \text{diag}(1, 0, 0, 0)$.

\[
    \mathbf{S}_A = \mathbf{W} \begin{pmatrix}1&0&0&0\\0&0&0&0\\0&0&0&0\\0&0&0&0\end{pmatrix} \mathbf{W}^{-1}
\]

First compute $\mathbf{W}\mathbf{D}_A$: this keeps only the first column of
$\mathbf{W}$ and zeros out the rest:
\[
    \mathbf{W}\mathbf{D}_A = \begin{pmatrix}1&0&0&0\\0&0&0&0\\1&0&0&0\\1&0&0&0\end{pmatrix}
\]

Now multiply by $\mathbf{W}^{-1}$ on the right. Since we only have non-zeros in
column 1 (corresponding to Gene A), and $[\mathbf{W}^{-1}]_{1,:} = [1,0,0,0]$
(first row of $\mathbf{W}^{-1}$):
\[
    \mathbf{S}_A = \begin{pmatrix}1&0&0&0\\0&0&0&0\\1&0&0&0\\1&0&0&0\end{pmatrix}
\]

\textbf{Reading $\mathbf{S}_A$:} Column $j$ of $\mathbf{S}_A$ is non-zero only
for $j = A$. Row $i$ is non-zero only for nodes $i$ that have Gene A as an
ancestor (i.e., $A, C, D$). This says: Gene A's causal channel is ``visible''
at nodes A, C, and D, but not at B (since A does not regulate B).
\end{examplebox}

\begin{examplebox}{Computing $\mathbf{S}_B$ (GSO for Gene B)}

$\text{pred}(B) = \{B\}$, so $\mathbf{D}_B = \text{diag}(0, 1, 0, 0)$.

By the same logic, $\mathbf{W}\mathbf{D}_B$ keeps only column 2 of $\mathbf{W}$:
\[
    \mathbf{W}\mathbf{D}_B = \begin{pmatrix}0&0&0&0\\0&1&0&0\\0&1&0&0\\0&1&0&0\end{pmatrix}
\]

Multiplying by $\mathbf{W}^{-1}$ and using $[\mathbf{W}^{-1}]_{2,:} = [0,1,0,0]$:
\[
    \mathbf{S}_B = \begin{pmatrix}0&0&0&0\\0&1&0&0\\0&1&0&0\\0&1&0&0\end{pmatrix}
\]

Gene B's causal channel is visible at B, C, and D. Not at A (B does not
regulate A).
\end{examplebox}

\begin{examplebox}{Computing $\mathbf{S}_C$ (GSO for Gene C)}

$\text{pred}(C) = \{A, B, C\}$, so $\mathbf{D}_C = \text{diag}(1, 1, 1, 0)$.

\[
    \mathbf{W}\mathbf{D}_C = \begin{pmatrix}1&0&0&0\\0&1&0&0\\1&1&1&0\\1&1&1&0\end{pmatrix}
\]

Now right-multiply by $\mathbf{W}^{-1}$:
\[
    \mathbf{S}_C = \mathbf{W}\mathbf{D}_C \mathbf{W}^{-1}
    = \begin{pmatrix}1&0&0&0\\0&1&0&0\\1&1&1&0\\1&1&1&0\end{pmatrix}
    \begin{pmatrix}1&0&0&0\\0&1&0&0\\-1&-1&1&0\\0&0&-1&1\end{pmatrix}
\]

Row 1 (Gene A): $[1, 0, 0, 0]$.
Row 2 (Gene B): $[0, 1, 0, 0]$.
Row 3 (Gene C): $[1\cdot1 + 1\cdot0 + 1\cdot(-1),\; 1\cdot0 + 1\cdot1 + 1\cdot(-1),\; 1\cdot0+1\cdot0+1\cdot1,\; 0]$
$= [0, 0, 1, 0]$.
Row 4 (Gene D): same as row 3 (since $\mathbf{W}\mathbf{D}_C$ has identical
rows 3 and 4): $[0, 0, 1, 0]$.

\[
    \mathbf{S}_C = \begin{pmatrix}1&0&0&0\\0&1&0&0\\0&0&1&0\\0&0&1&0\end{pmatrix}
\]

\textbf{Interpretation.} $[\mathbf{S}_C]_{D,C} = 1$: Gene D inherits Gene C's
causal channel. But the ``mixing'' through A and B cancels out, leaving a clean
projection onto Gene C's own contribution.
\end{examplebox}

\begin{examplebox}{Computing $\mathbf{S}_D$ (GSO for Gene D)}

$\text{pred}(D) = \{A, B, C, D\} = \mathcal{V}$, so $\mathbf{D}_D = \mathbf{I}$.

\[
    \mathbf{S}_D = \mathbf{W} \mathbf{I} \mathbf{W}^{-1} = \mathbf{W}\mathbf{W}^{-1} = \mathbf{I}
\]

The causal GSO for the terminal node is the identity matrix! This makes perfect
sense: Gene D is preceded by every gene in the network, so its ``causal
channel'' includes all frequencies. Applying $\mathbf{S}_D$ leaves the signal
unchanged.
\end{examplebox}

\paragraph{Summary of all four GSOs.}
\[
    \mathbf{S}_A = \begin{pmatrix}1&0&0&0\\0&0&0&0\\1&0&0&0\\1&0&0&0\end{pmatrix}, \quad
    \mathbf{S}_B = \begin{pmatrix}0&0&0&0\\0&1&0&0\\0&1&0&0\\0&1&0&0\end{pmatrix}
\]
\[
    \mathbf{S}_C = \begin{pmatrix}1&0&0&0\\0&1&0&0\\0&0&1&0\\0&0&1&0\end{pmatrix}, \quad
    \mathbf{S}_D = \begin{pmatrix}1&0&0&0\\0&1&0&0\\0&0&1&0\\0&0&0&1\end{pmatrix} = \mathbf{I}
\]

\subsection{Worked Numerical Example: Applying the GSOs to a Signal}

Let $\mathbf{x} = [2, 3, 1, 0.5]^\top$ (gene expression levels for A, B, C, D).

First, compute the exogenous contributions:
$\mathbf{c} = \mathbf{W}^{-1}\mathbf{x} = [2, 3, 1-2-3, 0.5-1]^\top = [2, 3, -4, -0.5]^\top$.

Now apply each GSO:

\begin{examplebox}{Numerical Application of Each Causal GSO}
\[
    \mathbf{S}_A \mathbf{x} = \begin{pmatrix}1\\0\\1\\1\end{pmatrix} \cdot x_A
    = \begin{pmatrix}2\\0\\2\\2\end{pmatrix}
\]
At every node that has A as an ancestor (A, C, D), the signal is replaced by
Gene A's own expression (2). Gene B, which does not inherit from A, gets 0.

\[
    \mathbf{S}_B \mathbf{x} = \begin{pmatrix}0\\1\\1\\1\end{pmatrix} \cdot x_B
    = \begin{pmatrix}0\\3\\3\\3\end{pmatrix}
\]

\[
    \mathbf{S}_C \mathbf{x} = \begin{pmatrix}x_A \\ x_B \\ x_C \\ x_C\end{pmatrix}
    = \begin{pmatrix}2\\3\\1\\1\end{pmatrix}
\]

\[
    \mathbf{S}_D \mathbf{x} = \mathbf{x} = \begin{pmatrix}2\\3\\1\\0.5\end{pmatrix}
\]

Observe that $\mathbf{S}_A\mathbf{x} + \mathbf{S}_B\mathbf{x} + \cdots$ does
\emph{not} simply reconstruct $\mathbf{x}$. The GSOs are not an orthonormal
decomposition; they are projection-like operators whose learnable combination
defines a graph filter.
\end{examplebox}

% =============================================================================
\section{Causal Graph Filters --- Linear Combinations of GSOs}
\label{sec:dagcn_filters}
% =============================================================================

With the causal GSOs defined, we can now construct proper \textbf{causal graph
filters}: linear combinations of the $N$ shift operators that perform a
well-defined ``convolution'' operation on DAG signals.

\subsection{Filter Definition}

\begin{mathbox}{Causal Graph Filter}
A causal graph filter $\mathbf{H}$ with parameters $\boldsymbol{\theta} =
[\theta_1, \dots, \theta_N]^\top \in \mathbb{R}^N$ is:
\begin{equation}
    \mathbf{H} = \sum_{k \in \mathcal{V}} \theta_k \mathbf{S}_k
    = \mathbf{W} \left(\sum_k \theta_k \mathbf{D}_k\right) \mathbf{W}^{-1}
    \label{eq:causal_filter}
\end{equation}

The diagonal matrix in the middle, $\boldsymbol{\Phi} = \sum_k \theta_k
\mathbf{D}_k$, is a \textbf{causal frequency response}: its $(i,i)$ entry is
$\Phi_{ii} = \sum_{k \ge i} \theta_k$ (the sum of filter coefficients for all
nodes that succeed $i$, including $i$ itself). In the DAG Fourier domain, the
filter multiplies the $i$-th causal frequency $c_i$ by
$\hat{h}_i = \sum_{k \ge i} \theta_k$.

\medskip
\textbf{Output.} For each node $i$:
\[
    [\mathbf{H}\mathbf{x}]_i = \sum_{j \le i} W_{ij} \, \hat{h}_j \, c_j
    = \sum_{j \le i} W_{ij} \left(\sum_{k \ge j} \theta_k\right) c_j
\]
Each ancestral exogenous contribution $c_j$ is re-weighted by the cumulative
filter coefficient $\hat{h}_j = \sum_{k \ge j} \theta_k$. This is a proper
convolution: it is shift-equivariant with respect to the causal GSOs and
reduces to standard polynomial filtering when the DAG is a directed path graph.
\end{mathbox}

\paragraph{Comparison to standard graph convolution.}
In a standard graph convolution, the filter $\mathbf{H} = \sum_r h_r
\mathbf{A}^r$ applies the same weight $h_r$ to contributions that are $r$ hops
away, regardless of direction. On a DAG, this means contributions from
descendants are treated identically to contributions from ancestors (once we
ignore directionality). By contrast, the causal filter $\mathbf{H} =
\sum_k \theta_k \mathbf{S}_k$ assigns a distinct learnable weight
$\theta_k$ to each node's causal channel, and \emph{only} aggregates from
ancestors (never from descendants). The filter is causal by construction.

\subsection{Worked Numerical Example}

\begin{examplebox}{Causal Filter with $\boldsymbol{\theta} = [1, 0.5, 0.5, 0.25]$ on $\mathbf{x} = [2, 3, 1, 0.5]$}

\textbf{Step 1.} Compute the frequency response at each node:
\begin{align*}
    \hat{h}_A &= \theta_A + \theta_B + \theta_C + \theta_D = 1 + 0.5 + 0.5 + 0.25 = 2.25 \\
    \hat{h}_B &= \theta_B + \theta_C + \theta_D = 0.5 + 0.5 + 0.25 = 1.25 \\
    \hat{h}_C &= \theta_C + \theta_D = 0.5 + 0.25 = 0.75 \\
    \hat{h}_D &= \theta_D = 0.25
\end{align*}

(Each $\hat{h}_j = \sum_{k \ge j} \theta_k$, summing over all descendants
of $j$ including $j$ itself.)

\textbf{Step 2.} Compute exogenous contributions:
$\mathbf{c} = [2, 3, -4, -0.5]^\top$ (computed earlier).

\textbf{Step 3.} Apply the filter. For each node $i$, sum $W_{ij} \hat{h}_j c_j$
over all ancestors $j \le i$:
\begin{align*}
    [\mathbf{H}\mathbf{x}]_A &= W_{AA} \hat{h}_A c_A
        = 1 \cdot 2.25 \cdot 2 = 4.5 \\
    [\mathbf{H}\mathbf{x}]_B &= W_{BB} \hat{h}_B c_B
        = 1 \cdot 1.25 \cdot 3 = 3.75 \\
    [\mathbf{H}\mathbf{x}]_C &= W_{CA} \hat{h}_A c_A
        + W_{CB} \hat{h}_B c_B
        + W_{CC} \hat{h}_C c_C \\
        &= 1 \cdot 2.25 \cdot 2 + 1 \cdot 1.25 \cdot 3 + 1 \cdot 0.75 \cdot (-4) \\
        &= 4.5 + 3.75 - 3 = 5.25 \\
    [\mathbf{H}\mathbf{x}]_D &= W_{DA}\hat{h}_A c_A + W_{DB}\hat{h}_B c_B
        + W_{DC}\hat{h}_C c_C + W_{DD}\hat{h}_D c_D \\
        &= 1 \cdot 2.25 \cdot 2 + 1 \cdot 1.25 \cdot 3 + 1 \cdot 0.75 \cdot (-4)
           + 1 \cdot 0.25 \cdot (-0.5) \\
        &= 4.5 + 3.75 - 3 - 0.125 = 5.125
\end{align*}

\textbf{Result:} $\mathbf{H}\mathbf{x} = [4.5, 3.75, 5.25, 5.125]^\top$.

Observe that no information flows from D back to A or B. The filter strictly
respects the causal structure of the DAG: the output at Gene A depends only on
Gene A's own exogenous contribution, and the output at Gene D aggregates all
four contributions weighted by their positions in the causal hierarchy.
\end{examplebox}

% =============================================================================
\section{The DCN Architecture}
\label{sec:dagcn_dcn}
% =============================================================================

With the causal graph filter in hand, Rey et al.\ construct a full deep learning
architecture called the \textbf{Directed Convolutional Network} (DCN). The DCN
is the multi-feature, multi-layer, nonlinear generalisation of the causal graph
filter.

\subsection{Multi-Feature Extension}

In practice, each node has not a single scalar signal but an
$F$-dimensional feature vector. We collect all node features into a matrix
$\mathbf{X}^{(\ell)} \in \mathbb{R}^{N \times F_\ell}$ where $N$ is the number
of nodes and $F_\ell$ is the feature dimension at layer $\ell$.

Each causal GSO $\mathbf{S}_k$ gets its own learnable weight matrix
$\boldsymbol{\Theta}_k^{(\ell)} \in \mathbb{R}^{F_\ell \times F_{\ell+1}}$ that
maps from the current feature dimension to the next. The DCN layer recursion is:

\begin{equation}
    \mathbf{X}^{(\ell+1)} = \sigma\!\left(
        \sum_{k \in \mathcal{V}} \mathbf{S}_k \, \mathbf{X}^{(\ell)} \,
        \boldsymbol{\Theta}_k^{(\ell)}
    \right)
    \label{eq:dcn_layer}
\end{equation}

where $\sigma = \text{ReLU}$ is the pointwise nonlinearity and $\mathbf{X}^{(0)}
= \mathbf{X}$ is the input feature matrix. Each $\mathbf{S}_k$ acts on all
$F_\ell$ input features simultaneously (it is applied to each column of
$\mathbf{X}^{(\ell)}$ independently), and the result is linearly mixed by
$\boldsymbol{\Theta}_k^{(\ell)}$ to produce the output features.

\subsection{Architecture Diagram}

\begin{figure}[H]
\centering
\begin{tikzpicture}[
    box/.style={rectangle, rounded corners=4pt, draw=purple!60!black,
                fill=purple!6!white, thick, minimum width=3.6cm,
                minimum height=0.85cm, font=\small\sffamily, align=center},
    gsobox/.style={rectangle, rounded corners=4pt, draw=orange!60!black,
                   fill=orange!5!white, thick, minimum width=3.6cm,
                   minimum height=0.85cm, font=\small\sffamily, align=center},
    arrow/.style={-{Stealth[length=3mm]}, thick, purple!60!black},
    label/.style={font=\scriptsize\sffamily\itshape, text=gray!60!black}
]
    \node[box]    (inp)  at (0,  0)   {Input $\mathbf{X}^{(0)} \in \mathbb{R}^{N \times F_0}$};
    \node[gsobox] (gso1) at (0, -1.8) {$\sum_{k} \mathbf{S}_k \mathbf{X}^{(0)} \boldsymbol{\Theta}_k^{(0)}$\\{\scriptsize $N$ causal GSO branches}};
    \node[box]    (act1) at (0, -3.4) {$\mathbf{X}^{(1)} = \sigma(\cdot) \in \mathbb{R}^{N \times F_1}$};
    \node[gsobox] (gso2) at (0, -5.2) {$\sum_{k} \mathbf{S}_k \mathbf{X}^{(1)} \boldsymbol{\Theta}_k^{(1)}$\\{\scriptsize same structure, new weights}};
    \node[box]    (act2) at (0, -6.8) {$\mathbf{X}^{(2)} = \sigma(\cdot) \in \mathbb{R}^{N \times F_2}$};
    \node[box]    (out)  at (0, -8.4) {Task head (classification / regression)};

    \draw[arrow] (inp)  -- (gso1);
    \draw[arrow] (gso1) -- (act1);
    \draw[arrow] (act1) -- (gso2);
    \draw[arrow] (gso2) -- (act2);
    \draw[arrow] (act2) -- node[label, right]{final layer} (out);

    % Side annotation
    \node[label, align=left] at (5.0, -1.8)
        {$N$ weight matrices\\$\boldsymbol{\Theta}_k^{(0)} \in \mathbb{R}^{F_0 \times F_1}$};
    \node[label, align=left] at (5.0, -5.2)
        {$N$ weight matrices\\$\boldsymbol{\Theta}_k^{(1)} \in \mathbb{R}^{F_1 \times F_2}$};
\end{tikzpicture}
\caption{The \textbf{Directed Convolutional Network (DCN)} architecture. At each
layer, all $N$ causal GSOs are applied in parallel to the current feature
matrix, each with its own weight matrix. The results are summed and passed
through a pointwise nonlinearity. The total parameter count at layer $\ell$ is
$N \times F_\ell \times F_{\ell+1}$ --- it \emph{scales with the graph size}
$N$.}
\label{fig:dcn_architecture}
\end{figure}

\subsection{Forward Pass: Dimensions and Computation}

Let us trace through the DCN forward pass with concrete dimensions for the gene
network ($N=4$ nodes, $F_0=2$ input features, $F_1=4$ hidden features,
$F_2=1$ output).

\begin{examplebox}{Forward Pass Walkthrough for the Gene Network}

\textbf{Input.} $\mathbf{X}^{(0)} \in \mathbb{R}^{4 \times 2}$: two features
per gene (e.g., expression level and sequence motif score).

\textbf{Layer 0.} For each of the 4 nodes $k \in \{A, B, C, D\}$:
\begin{enumerate}[noitemsep]
    \item Compute $\mathbf{S}_k \mathbf{X}^{(0)} \in \mathbb{R}^{4 \times 2}$
          (apply the $4 \times 4$ causal GSO to the $4 \times 2$ feature matrix).
    \item Multiply by $\boldsymbol{\Theta}_k^{(0)} \in \mathbb{R}^{2 \times 4}$
          to get a $4 \times 4$ contribution.
\end{enumerate}
Sum all 4 contributions: $\sum_k \mathbf{S}_k \mathbf{X}^{(0)} \boldsymbol{\Theta}_k^{(0)}
\in \mathbb{R}^{4 \times 4}$.

Apply ReLU: $\mathbf{X}^{(1)} = \text{ReLU}(\cdot) \in \mathbb{R}^{4 \times 4}$.

\textbf{Layer 1.} Repeat: for each of 4 GSOs, apply to $\mathbf{X}^{(1)} \in
\mathbb{R}^{4 \times 4}$ and multiply by $\boldsymbol{\Theta}_k^{(1)} \in
\mathbb{R}^{4 \times 1}$. Sum, apply ReLU:
$\mathbf{X}^{(2)} \in \mathbb{R}^{4 \times 1}$.

\textbf{Parameter count.} Layer 0: $4 \times 2 \times 4 = 32$ parameters.
Layer 1: $4 \times 4 \times 1 = 16$ parameters. Total: 48 parameters (plus
task head). For a graph of $N$ nodes with $F$ features and $L$ layers, the
total is $O(N F^2 L)$ --- \emph{linear in} $N$.
\end{examplebox}

\subsection{Permutation Equivariance}

The DCN satisfies \textbf{permutation equivariance} (Corollary~1 in the paper):
if we relabel the nodes of the DAG by any permutation $\pi$ that preserves the
DAG structure (i.e., an automorphism), the output is permuted by the same
$\pi$. More precisely, for any topological-order-preserving permutation $\pi$:
\[
    \mathbf{H}(\mathbf{P}_\pi \mathbf{X}) = \mathbf{P}_\pi \mathbf{H}(\mathbf{X})
\]
where $\mathbf{P}_\pi$ is the permutation matrix corresponding to $\pi$. This
is the fundamental equivariance requirement for any respectable graph learning
architecture.

\subsection{DCN-T: The Transposed Variant for Source Detection}

For source detection tasks, information flows in the \emph{opposite} direction
to the DAG's causal edges. We observe who is infected (downstream) and want to
identify the source (upstream). To support backward-looking aggregation, Rey et
al.\ define the \textbf{transposed DCN} (DCN-T) by replacing each causal GSO
$\mathbf{S}_k$ with its transpose $\mathbf{S}_k^\top$:

\begin{equation}
    \mathbf{X}^{(\ell+1)}_{\text{DCN-T}} = \sigma\!\left(
        \sum_{k \in \mathcal{V}} \mathbf{S}_k^\top \, \mathbf{X}^{(\ell)} \,
        \boldsymbol{\Theta}_k^{(\ell)}
    \right)
    \label{eq:dcnt_layer}
\end{equation}

The transposed GSO $\mathbf{S}_k^\top$ aggregates information from a node's
\emph{descendants} rather than its ancestors. For source detection: when we
apply DCN-T to the infection state, each candidate source node's representation
accumulates evidence from all nodes it could have infected (its descendants).
The true source will have the richest set of infected descendants, and DCN-T
is designed to detect exactly this pattern.

\subsection{Computational Complexity}

At each DCN layer, the dominant cost is the $N$ matrix products
$\mathbf{S}_k \mathbf{X}^{(\ell)}$. Each $\mathbf{S}_k$ has at most
$\|\mathbf{S}_k\|_0 \le N^2$ non-zero entries (in practice much fewer for
sparse DAGs). Let $S = \max_k \|\mathbf{S}_k\|_0$ be the maximum sparsity.
The total cost per layer is $O(N \cdot S \cdot F_\ell)$ for the GSO
applications, plus $O(N \cdot F_\ell \cdot F_{\ell+1})$ for the weight
multiplications. On sparse DAGs (e.g., trees or chains), $S = O(N)$, giving
an overall per-layer cost of $O(N^2 F + NF^2)$.

% =============================================================================
\section{Parallel DCN (PDCN) --- Trading Depth for Width}
\label{sec:dagcn_pdcn}
% =============================================================================

The DCN architecture as defined in Equation~(\ref{eq:dcn_layer}) has one
important limitation: its parameter count scales with $N$ (the graph size),
because there are $N$ separate weight matrices $\boldsymbol{\Theta}_k^{(\ell)}$
per layer. For large graphs, this becomes prohibitive. Moreover, the authors
show that the expressiveness gain from stacking multiple DCN layers is limited
by a composition lemma.

\subsection{The Composition Lemma: Depth Does Not Help}

\begin{conceptbox}{Why Deep DCN Has Limited Expressiveness}
\textbf{Lemma (Composition of causal filters).}
The composition of two causal graph filters is itself a causal graph filter.
Formally, if $\mathbf{H}_1 = \mathbf{W}\boldsymbol{\Phi}_1\mathbf{W}^{-1}$ and
$\mathbf{H}_2 = \mathbf{W}\boldsymbol{\Phi}_2\mathbf{W}^{-1}$, then:
\[
    \mathbf{H}_1 \mathbf{H}_2 = \mathbf{W}(\boldsymbol{\Phi}_1\boldsymbol{\Phi}_2)\mathbf{W}^{-1}
\]
which is another causal graph filter with frequency response
$\boldsymbol{\Phi}_1\boldsymbol{\Phi}_2$ (a diagonal matrix, since the product
of two diagonal matrices is diagonal).

\medskip
\textbf{Consequence.} Two DCN layers (ignoring the nonlinearity) collapse into
one DCN layer with a modified frequency response. Adding more layers does not
introduce new representational capabilities that could not be achieved by a
single deeper filter. In the absence of the nonlinearity, a deep linear DCN is
equivalent to a single causal graph filter. The nonlinearity $\sigma$ breaks
this collapse, but only partially: empirically, deep DCN does not offer
substantial gains over a single layer followed by a nonlinearity, and adding
layers primarily increases computational cost without a commensurate accuracy
improvement.
\end{conceptbox}

This motivates a different scaling strategy: rather than going \emph{deep}
(stacking layers), go \emph{wide} (running $N$ parallel branches with shared
weights).

\subsection{The PDCN Architecture}

\begin{mathbox}{Parallel DCN (PDCN)}
The PDCN architecture has $N$ parallel processing branches, one for each causal
GSO. All branches share the same MLP parameters $\boldsymbol{\Theta}^{(\ell)}$:

\begin{equation}
    \mathbf{Z}_k^{(\ell)} = \text{MLP}\!\left(\mathbf{S}_k \mathbf{X}^{(\ell)};
    \boldsymbol{\Theta}^{(\ell)}\right), \quad k \in \mathcal{V}
    \label{eq:pdcn_branch}
\end{equation}

\begin{equation}
    \mathbf{X}^{(\ell+1)} = \sum_{k \in \mathcal{V}} \mathbf{Z}_k^{(\ell)}
    \label{eq:pdcn_sum}
\end{equation}

Each branch applies the same MLP (with ReLU activations and $L$ layers of
hidden dimension $F$) to a different GSO-transformed version of the feature
matrix. The results are summed across all $N$ branches.
\end{mathbox}

\subsection{Architecture Diagram: DCN vs. PDCN}

\begin{figure}[H]
\centering
\begin{tikzpicture}[
    inp/.style={rectangle, rounded corners=3pt, draw=blue!60!black,
                fill=blue!5!white, thick, minimum width=2.6cm,
                minimum height=0.75cm, font=\small\sffamily, align=center},
    branch/.style={rectangle, rounded corners=3pt, draw=orange!60!black,
                   fill=orange!5!white, thick, minimum width=2.1cm,
                   minimum height=0.75cm, font=\small\sffamily, align=center},
    sumbox/.style={circle, draw=red!60!black, fill=red!5!white, thick,
                   minimum size=0.75cm, font=\bfseries\sffamily},
    arrow/.style={-{Stealth[length=2.5mm]}, thick},
    label/.style={font=\scriptsize\sffamily\itshape, text=gray!60!black}
]
    % Input
    \node[inp] (X) at (0, 0) {$\mathbf{X}^{(\ell)}$};

    % Four branches (one per GSO)
    \node[branch] (B1) at (-4.5, -1.8) {$\mathbf{S}_A\mathbf{X}^{(\ell)}$};
    \node[branch] (B2) at (-1.5, -1.8) {$\mathbf{S}_B\mathbf{X}^{(\ell)}$};
    \node[branch] (B3) at ( 1.5, -1.8) {$\mathbf{S}_C\mathbf{X}^{(\ell)}$};
    \node[branch] (B4) at ( 4.5, -1.8) {$\mathbf{S}_D\mathbf{X}^{(\ell)}$};

    % MLP boxes (shared weights)
    \node[branch] (M1) at (-4.5, -3.4) {\small MLP($\cdot$; $\boldsymbol{\Theta}$)};
    \node[branch] (M2) at (-1.5, -3.4) {\small MLP($\cdot$; $\boldsymbol{\Theta}$)};
    \node[branch] (M3) at ( 1.5, -3.4) {\small MLP($\cdot$; $\boldsymbol{\Theta}$)};
    \node[branch] (M4) at ( 4.5, -3.4) {\small MLP($\cdot$; $\boldsymbol{\Theta}$)};

    % Sum node
    \node[sumbox] (S) at (0, -5.2) {$+$};

    % Output
    \node[inp] (Xp) at (0, -6.8) {$\mathbf{X}^{(\ell+1)}$};

    % Arrows: input to branches
    \draw[arrow] (X) -- (B1);
    \draw[arrow] (X) -- (B2);
    \draw[arrow] (X) -- (B3);
    \draw[arrow] (X) -- (B4);

    % Arrows: branches to MLPs
    \draw[arrow] (B1) -- (M1);
    \draw[arrow] (B2) -- (M2);
    \draw[arrow] (B3) -- (M3);
    \draw[arrow] (B4) -- (M4);

    % Arrows: MLPs to sum
    \draw[arrow] (M1) -- (S);
    \draw[arrow] (M2) -- (S);
    \draw[arrow] (M3) -- (S);
    \draw[arrow] (M4) -- (S);

    % Sum to output
    \draw[arrow] (S) -- (Xp);

    % Shared weights label
    \node[label] at (6.5, -3.4) {Same $\boldsymbol{\Theta}$\\for all branches};
    \draw[->, gray!40, thin] (5.8, -3.4) -- (5.0, -3.4);
\end{tikzpicture}
\caption{The \textbf{Parallel DCN (PDCN)} architecture for the gene network
($N = 4$). The input $\mathbf{X}^{(\ell)}$ is first transformed by each of the
$N = 4$ causal GSOs in parallel. All $N$ resulting matrices are processed by
the \emph{same} shared MLP (weight-tied across branches), then summed to
produce $\mathbf{X}^{(\ell+1)}$. Parameter count: $O(F^2 L)$ independently
of $N$.}
\label{fig:pdcn_architecture}
\end{figure}

\subsection{Parameter Count: Independent of Graph Size}

The decisive advantage of the PDCN over the DCN is parameter efficiency. In the
DCN, there are $N$ separate weight matrices per layer, giving $O(NF^2L)$ total
parameters. In the PDCN, all $N$ branches share a single MLP, so the total
parameter count is $O(F^2 L)$ --- completely \emph{independent of} $N$.

This has three consequences:

\begin{enumerate}
    \item \textbf{Scalability.} The PDCN can be applied to large DAGs without
          increasing the number of learned parameters. The same model can handle
          DAGs with 10 nodes or 10,000 nodes.
    \item \textbf{Generalisation.} Because weights are shared across branches,
          the PDCN is forced to learn features that are meaningful for any
          GSO-transformed input, not just the one corresponding to a specific
          node $k$. This acts as a form of regularisation.
    \item \textbf{Permutation equivariance (Remark~1).} Sharing weights across
          branches ensures that the PDCN is permutation equivariant: relabelling
          the nodes produces the same output, relabelled identically.
\end{enumerate}

\begin{keyinsight}{PDCN Achieves WL-Expressiveness}
The PDCN can distinguish any pair of non-isomorphic DAGs that differ in their
causal GSO structure. Since the causal GSOs $\{\mathbf{S}_k\}_{k \in
\mathcal{V}}$ uniquely characterise the DAG up to isomorphism (Proposition~2 in
the paper), any function of the full set of GSO-transformed features
$\{\mathbf{S}_k \mathbf{X}\}_{k \in \mathcal{V}}$ is at least as expressive as
the Weisfeiler--Lehman (WL) graph isomorphism test. The PDCN, by processing
each GSO branch with a shared MLP and summing, achieves exactly this level of
expressiveness.
\end{keyinsight}

\subsection{When to Use DCN vs. PDCN}

The choice between DCN and PDCN depends on the specific requirements:

\begin{center}
\begin{tabular}{lll}
\toprule
\textbf{Criterion} & \textbf{DCN} & \textbf{PDCN} \\
\midrule
Parameter count & $O(NF^2L)$ --- scales with $N$ & $O(F^2L)$ --- independent of $N$ \\
Expressiveness & High (per-node weights) & High (WL-level, shared weights) \\
Inductive learning & No (tied to graph size) & Yes (same model, any size) \\
When $N$ is large & Expensive & Preferred \\
When $N$ is small & Both applicable & Preferred for generalisation \\
Source detection & DCN-T (transposed) & PDCN-T (transposed) \\
\bottomrule
\end{tabular}
\end{center}

% =============================================================================
\section{Source Detection with DCN-T}
\label{sec:dagcn_source}
% =============================================================================

The most compelling application of the DAGCN framework in the paper is
\textbf{source identification}: given a signal propagated on a DAG (for example,
a viral cascade on a gene regulatory network), identify the node that initiated
the propagation.

\subsection{The Source Identification Task}

In the source identification setting, we observe the final signal state
$\mathbf{x}_{\text{obs}} \in \mathbb{R}^N$ at every node after the propagation
has run to completion. Under the SEM, this signal was generated as $\mathbf{x}
= \mathbf{W}\mathbf{c}$ where $\mathbf{c}$ has exactly one large non-zero entry
at the source node $k^*$ (and small noise contributions at other nodes). The
task is to identify $k^*$ from $\mathbf{x}_{\text{obs}}$.

\begin{examplebox}{Why This Is Hard}
Suppose the true source is Gene A in our network. Gene A's exogenous
contribution $c_A$ propagates to all its descendants: C and D each accumulate
a contribution of $W_{CA} c_A = c_A$ and $W_{DA} c_A = c_A$ respectively.
The signal at Gene D is $x_D = c_A + c_B + c_C + c_D \approx c_A$ if $c_A$
dominates. But from $x_D$ alone, we cannot tell whether the large signal came
from Gene A or Gene B (both equally contribute to Gene D). We need to look at
the \emph{pattern} across all nodes to identify the source.

In particular, if Gene B's contribution is small ($c_B \approx 0$) but Gene A
is the source ($c_A$ large), then: $x_A = c_A$ (large), $x_B \approx 0$ (small),
$x_C \approx c_A$ (large), $x_D \approx c_A$ (large). The signature is:
large signal at A, C, D; small at B. A model looking only at local node
features or only at direct edges would miss this pattern; it requires reasoning
about the \emph{transitive} influence of Gene A.
\end{examplebox}

\subsection{Why Transposed GSOs Are Needed}

The standard causal GSO $\mathbf{S}_k$ aggregates from a node's \emph{ancestors}
(looking upstream). For source detection, we need the opposite: we need each
candidate source node to aggregate from its \emph{descendants} (looking
downstream). The reason is intuitive --- we want the source node to ``see''
how many and which downstream nodes exhibit the propagated signal.

The transposed GSO $\mathbf{S}_k^\top$ does exactly this. The $(i,j)$ entry of
$\mathbf{S}_k^\top$ is $[\mathbf{S}_k]_{ji}$, which is non-zero when node $i$
is an ancestor of node $j$ and both $j$ and $k$ share a common ancestor. In the
context of source detection, applying $\mathbf{S}_k^\top$ to the observed signal
allows candidate source node $k$ to collect evidence from all nodes downstream
of it:

\[
    [\mathbf{S}_k^\top \mathbf{x}]_i = \sum_{j: i \le j,\; k \le j}
    [\mathbf{W}^\top]_{ij} \, [\mathbf{W}^{-\top}]_{ji} \, x_j
\]

In plain terms: $\mathbf{S}_k^\top$ at node $i$ aggregates the observed signals
of all nodes $j$ that are descendants of both $i$ and $k$. If $i = k^*$ (the
true source), this aggregates evidence from every infected descendant.

\begin{analogybox}{Finding the Source of a Rumour by Asking Who Heard It}
Imagine a rumour that started with one person and spread through a social
network. Each person only passes the rumour to people they directly communicate
with (causal direction: source $\to$ recipients). After the rumour has spread,
you want to find who started it.

You cannot ask who \emph{told} you the rumour (that would be aggregating from
ancestors, looking upstream --- the DCN direction). Instead, you ask each person:
``Of all the people who know the rumour, which ones could only have heard it from
you?'' That question looks \emph{downstream}: it asks who is in your rumour
propagation cone. The person whose downstream cone matches the set of informed
people most precisely is the source. This is exactly what $\mathbf{S}_k^\top$
computes for each candidate source node $k$.
\end{analogybox}

\subsection{Connection to Reversed DAG Approaches}

The DCN-T approach is conceptually aligned with the reversed-DAG message
passing used in the Saram\"{a}ki Temporal Event Graph framework (Chapter~\ref{chap:saramaki}).
Both recognise the same fundamental insight: epidemic source detection is a
\emph{backward inference} problem. The observed infection pattern is the
downstream consequence of the source's action; we must propagate evidence
upstream (in the Saram\"{a}ki approach, by flipping the DAG edges) or equivalently
aggregate from descendants (in the DCN-T approach, by transposing the GSOs).

The key difference is in the representation: the Saram\"{a}ki approach builds
an explicit event-level DAG and reverses it structurally, while the DCN-T works
on the original node-level DAG and reverses the direction of aggregation through
matrix transposition. The underlying causal logic is identical.

\subsection{Experimental Results}

Rey et al.\ evaluate source identification accuracy on synthetic DAGs generated
from Erd\H{o}s--R\'{e}nyi and Barab\'{a}si--Albert DAG models. The source node
generates a large exogenous input $c_{k^*} \gg 0$ and all other nodes have
small noise contributions. The models predict the most likely source node.

The results are striking:

\begin{center}
\begin{tabular}{lcc}
\toprule
\textbf{Model} & \textbf{Accuracy} & \textbf{Notes} \\
\midrule
Standard DCN (forward GSOs) & 0.057 & Near-random; cannot look downstream \\
DCN-T (transposed GSOs) & \textbf{0.997} & Near-perfect; correctly aggregates downstream \\
Standard GCN (undirected) & 0.312 & Partial success; loses causal direction \\
\bottomrule
\end{tabular}
\end{center}

The contrast between standard DCN (0.057 accuracy, essentially random) and
DCN-T (0.997 accuracy, near-perfect) is perhaps the most informative result
in the paper. It is not that the architecture is wrong or the data is
insufficient --- the standard DCN is simply looking in the wrong direction. By
transposing the GSOs, the DCN-T aggregates exactly the right information, and
the task becomes almost trivially easy. The causal direction of aggregation is
not a hyperparameter to be tuned; it is a structural requirement of the task.

The standard GCN achieves 0.312 accuracy --- well above random (0.25 for a
four-class problem) but far below DCN-T. This is because undirected aggregation
captures some of the signal (a node with many infected neighbours is likely
near the source) but corrupts the representation by mixing upstream and
downstream information symmetrically.

% =============================================================================
\section{Summary}
\label{sec:dagcn_summary}
% =============================================================================

This chapter has developed the complete theoretical and practical framework
of Rey, Ajorlou, and Mateos's \textbf{Directed Acyclic Graph Convolutional
Network (DAGCN)}~\cite{rey2025dagcn}. We close with a structured comparison
of the three architectures and a discussion of the paper's position in the
broader landscape.

\subsection{Comparison Table}

\begin{center}
\begin{tabular}{llll}
\toprule
\textbf{Property} & \textbf{Standard GNN} & \textbf{DCN} & \textbf{PDCN} \\
\midrule
Handles DAG structure & No & Yes & Yes \\
Spectral foundation & Eigendecomposition & Transitive closure & Transitive closure \\
Aggregation direction & Both / undirected & Ancestors only & Ancestors only \\
Basis & Eigenvectors of $\mathbf{L}$ & Columns of $\mathbf{W}$ & Columns of $\mathbf{W}$ \\
Shift operators & Powers of $\mathbf{A}$ or $\mathbf{L}$ & $\{\mathbf{S}_k\}_{k \in \mathcal{V}}$ & $\{\mathbf{S}_k\}_{k \in \mathcal{V}}$ \\
Param count / layer & $O(RF^2)$ & $O(NF^2)$ & $O(F^2L)$ \\
Scales with graph & No & Yes (bad) & No (good) \\
Source detection & Poor & Poor (forward) & Excellent (transposed) \\
WL expressiveness & Limited & Yes & Yes \\
Inductive (new graphs) & Limited & No (fixed $N$) & Yes \\
\bottomrule
\end{tabular}
\end{center}

\subsection{The Three Key Contributions}

Rey et al.\ make three tightly coupled contributions:

\begin{enumerate}
    \item \textbf{The causal GSO and DAG Fourier theory.} By recognising the
          transitive closure $\mathbf{W} = (\mathbf{I} - \mathbf{A})^{-1}$ as the
          natural Fourier basis for DAG signals and defining the causal GSOs
          $\mathbf{S}_k = \mathbf{W}\mathbf{D}_k\mathbf{W}^{-1}$, the paper
          provides a rigorous spectral theory for DAGs that avoids the nilpotency
          problem entirely. The spectrum is no longer the (degenerate) eigenvalues
          of $\mathbf{A}$; it is the per-node causal channels, one for each node
          in the graph.

    \item \textbf{The DCN and PDCN architectures.} The DCN generalises the
          causal filter to a deep, multi-feature, nonlinear architecture with
          per-node weight matrices. The PDCN resolves the scaling problem by
          sharing weights across GSO branches, achieving $O(F^2L)$ parameter
          count regardless of graph size. Both architectures are permutation
          equivariant and WL-expressive.

    \item \textbf{The DCN-T transposition for causal inference.} The insight
          that source detection requires \emph{transposed} GSOs --- aggregating
          from descendants rather than ancestors --- is both elegant and
          empirically validated. The accuracy gap between DCN (0.057) and DCN-T
          (0.997) demonstrates that architectural alignment with the task's
          causal direction is not a minor improvement but a categorical difference.
\end{enumerate}

\subsection{Position in the Broader Landscape}

The DAGCN paper occupies a distinctive position in the GNN landscape for
temporal and causal networks. Compared to the approaches studied in previous
chapters:

\paragraph{Versus Saram\"{a}ki's TEG (Chapter~\ref{chap:saramaki}).} Both
approaches recognise that DAG structure is fundamental to causal reasoning on
temporal networks. The TEG approach lifts the problem to the event level,
constructing a new DAG of events and running GNN message passing on it. The
DAGCN approach works directly on the original node-level DAG, using the
transitive closure to build a rich spectral theory. The TEG approach handles
arbitrary temporal networks (including those with undirected contacts), while
DAGCN requires a pre-specified DAG structure.

\paragraph{Versus standard temporal GNNs.} Temporal GNNs such as TGAT or
TGN use attention mechanisms or memory modules to order temporal information,
but they operate on the original node graph and do not exploit DAG structure
explicitly. DAGCN, by contrast, uses the full ancestry structure encoded in
$\mathbf{W}$ to define shift operators. On true DAG data, DAGCN's structural
advantage translates directly into better empirical performance.

\paragraph{Versus backtracking-based methods.} Methods such as the
Backtracking Network of Ru et al.\ use explicit probabilistic likelihood models
to score candidate sources. DAGCN learns these scores end-to-end from data,
without requiring a parametric model of the propagation process. This makes
DAGCN more flexible in settings where the generative model is unknown, at the
cost of requiring labelled training examples of (source, observed signal) pairs.

\begin{keyinsight}{The Central Lesson of This Chapter}
The central lesson of the DAGCN paper is that the appropriate notion of
``graph frequency'' depends fundamentally on the type of graph being processed.
For undirected graphs, the eigenvectors of the Laplacian are the right basis.
For DAGs, the eigenvectors of $\mathbf{A}$ are degenerate and useless; the
right basis is the columns of the transitive closure $\mathbf{W}$. A GNN
designed for undirected graphs cannot simply be applied to a DAG by ignoring
the direction --- it will be blind to the very structure that makes the DAG
informative. The causal GSOs $\{\mathbf{S}_k\}$ provide the first principled,
expressive, and computationally tractable signal processing framework for DAGs,
and the DCN/PDCN architectures operationalise this framework in a learnable,
scalable form.
\end{keyinsight}
